\documentclass[a4paper,11pt]{article}
\usepackage[utf8]{inputenc}
\usepackage[T1]{fontenc}
\usepackage[ngerman]{babel}
\usepackage{geometry}
\geometry{left=2.5cm, right=2.5cm, top=2.5cm, bottom=2.5cm}
\usepackage{longtable}
\usepackage{booktabs}
\usepackage{array}
\usepackage{xcolor}
\usepackage{colortbl}
\usepackage{hyperref}
\usepackage{fancyhdr}
\usepackage{tabularx}
\usepackage{listings}
\usepackage{enumitem}
\usepackage{graphicx}

\definecolor{headerblue}{HTML}{1565C0}
\definecolor{tableheader}{HTML}{E3F2FD}
\definecolor{codebg}{HTML}{F5F5F5}
\definecolor{codegreen}{HTML}{2E7D32}
\definecolor{codegray}{HTML}{757575}

\lstset{
  backgroundcolor=\color{codebg},
  basicstyle=\ttfamily\small,
  keywordstyle=\color{headerblue}\bfseries,
  commentstyle=\color{codegreen},
  stringstyle=\color{codegray},
  breaklines=true,
  frame=single,
  rulecolor=\color{codegray!30},
  numbers=none,
  showstringspaces=false,
  tabsize=2
}

\hypersetup{
    colorlinks=true,
    linkcolor=headerblue,
    urlcolor=headerblue
}

\pagestyle{fancy}
\fancyhf{}
\fancyhead[L]{\textbf{INSPECTAIR}}
\fancyhead[R]{Design Description v2.0}
\fancyfoot[C]{\thepage}
\renewcommand{\headrulewidth}{0.4pt}

\setlength{\parindent}{0pt}
\setlength{\parskip}{6pt}

\begin{document}

\begin{center}
{\LARGE\textbf{INSPECTAIR}}\\[0.3cm]
{\Large Design Description}\\[0.3cm]
{\large System- und Software-Architektur}\\[0.5cm]
\begin{tabular}{ll}
\textbf{Projekt:} & InspectAir -- Luftqualitätsmessgerät \\
\textbf{Version:} & 2.0 \\
\textbf{Datum:} & Februar 2026 \\
\textbf{Autoren:} & Maximilian Liebl, Cedric Stroh, Jannis Messer \\
\textbf{Kurs:} & Mikrocomputertechnik 1, DHBW Ravensburg \\
\textbf{Betreuer:} & Thomas Stingl (Airbus) \\
\end{tabular}
\end{center}

\vspace{0.5cm}
\noindent\textit{Dieses Dokument wurde unter Verwendung von KI-Tools (Claude, Anthropic) zur Überprüfung erstellt.}
\vspace{0.3cm}

\tableofcontents
\newpage

% ============================================================
\section{Projektbeschreibung}
% ============================================================

\subsection{Überblick}
InspectAir ist ein Luftqualitätsmonitor für Innenräume. Das System erfasst kontinuierlich relevante Luftqualitätsparameter und stellt diese auf einem 4-Zoll-TFT-Display dar. Durch Farbcodierung können kritische Werte auf einen Blick erkannt werden.

\subsection{Kernfunktionen}
\begin{itemize}[nosep]
\item Messung von CO\textsubscript{2}, Feinstaub (PM1.0/PM2.5/PM10), VOC, Temperatur und Luftfeuchtigkeit
\item Echtzeit-Anzeige aller Messwerte auf 4-Zoll-Display (480$\times$320 Pixel)
\item Vierstufige Ampel-Farbcodierung (Grün/Gelb/Orange/Rot) nach WHO/UBA-Grenzwerten
\item Fünf wählbare UI-Designs (Tree, Overview, Detail, Analog, Bubble)
\item Präsenzerkennung via 24\,GHz FMCW-Radar für automatisches Display-Management
\item Energiesparmodi: Display-Dimming und Light Sleep
\end{itemize}

\subsection{Zielgruppe}
Das Gerät richtet sich an Privatpersonen und Büros, die ihre Raumluftqualität überwachen möchten. Typische Einsatzorte sind Wohnzimmer, Schlafzimmer, Büros und Klassenzimmer.

% ============================================================
\section{System-Architektur}
% ============================================================

\subsection{Hardware-Blockdiagramm}

Das System besteht aus folgenden Hardware-Blöcken:

\begin{lstlisting}[language={},basicstyle=\ttfamily\footnotesize]
+---------------------------------------------------------------+
|                    INSPECTAIR SYSTEM                           |
+---------------------------------------------------------------+
|                                                               |
|  +--------+  +--------+  +--------+  +--------+              |
|  | AHT20  |  | SGP40  |  | MH-Z19C|  | PMS5003|              |
|  |Temp/Hum|  |  VOC   |  |  CO2   |  |Feinst. |              |
|  +---+----+  +---+----+  +---+----+  +---+----+              |
|      |I2C        |I2C        |SoftSer    |UART1               |
|      +-----+-----+          |           |                     |
|            |                 |           |                     |
|  +---------+--------+-------+-----------+--------+            |
|  |                                               |            |
|  |          ESP32-S3-WROOM N8R8                   |            |
|  |             (Waveshare)                        |            |
|  |                                               |            |
|  +--------+------------------+------ ---+--------+            |
|           |SPI               |GPIO       |UART2               |
|    +------+------+    +------+------+  +-+--------+           |
|    |  ST7796S    |    |   BS170     |  | LD2410C  |           |
|    |  4" TFT     |    |   MOSFET    |  |  Radar   |           |
|    |  Display    |    |  Backlight  |  | (Level   |           |
|    +-------------+    +-------------+  | Shifter) |           |
|                                        +----------+           |
+---------------------------------------------------------------+
\end{lstlisting}

\subsection{Schnittstellen-Übersicht}

\begin{longtable}{p{2.2cm} p{2cm} p{1.5cm} p{3cm} p{3.5cm}}
\toprule
\rowcolor{tableheader}
\textbf{Komponente} & \textbf{Interface} & \textbf{Spannung} & \textbf{ESP32 Pins} & \textbf{Besonderheit} \\
\midrule
\endhead

ST7796S Display & SPI & 3,3\,V & 11(MOSI), 12(SCLK), 14(CS), 9(DC), 46(RST), 3(BL) & 480$\times$320 Pixel \\
\midrule
AHT20 & I\textsuperscript{2}C (0x38) & 3,3\,V & 8 (SDA), 18 (SCL) & Pull-ups 4,7\,k$\Omega$ \\
\midrule
SGP40 & I\textsuperscript{2}C (0x59) & 3,3\,V & 8 (SDA), 18 (SCL) & Parallel zu AHT20 \\
\midrule
MH-Z19C & SoftSerial 9600 & 5\,V & 4 (RX), 5 (TX) & 220\,µF Stützkond.! \\
\midrule
PMS5003 & UART1 9600 & 5\,V & 16 (RX), 17 (TX) & Molex-Stecker \\
\midrule
LD2410C & UART2 256000 & 5\,V & 7 (RX), 6 (TX), 15 (OUT) & Level Shifter + 220\,µF! \\
\midrule
BS170 MOSFET & PWM & 3,3\,V Gate & 3 & 1\,k$\Omega$ Gate-Widerstand \\
\midrule
Button & GPIO & 3,3\,V & 1 & Interner Pull-up \\

\bottomrule
\end{longtable}

\subsection{Kritische Hardware-Design-Entscheidungen}

\subsubsection{Stützkondensatoren (220\,µF 25\,V)}
\textbf{Problem:} MH-Z19C und LD2410C verursachen Stromspitzen beim Messvorgang, die zu Spannungseinbrüchen und Systemabstürzen führen.

\textbf{Lösung:} 220\,µF Elkos (25\,V) direkt an VCC/GND beider Sensoren puffern die Stromspitzen ab.

\subsubsection{Level Shifter für LD2410C}
\textbf{Problem:} Der LD2410C sendet TX-Signale mit 5\,V-Pegel, die ESP32-GPIOs sind nur 3,3\,V-tolerant.

\textbf{Lösung:} Bidirektionaler Level Shifter (TXS0108E) zwischen Radar-TX und ESP32-RX. Die RX-Richtung (ESP$\rightarrow$Radar) kann direkt erfolgen, da der Radar 3,3\,V akzeptiert.

\subsubsection{BS170 MOSFET für Backlight}
\textbf{Problem:} Das Display-Backlight benötigt mehr Strom als ein GPIO liefern kann, und PWM-Dimming ist gewünscht.

\textbf{Lösung:} BS170 N-Kanal MOSFET als Low-Side-Switch. Logic-Level-MOSFET schaltet bei 3,3\,V durch. Gate über 1\,k$\Omega$ an GPIO3 für PWM-Steuerung.

\subsubsection{SoftwareSerial für MH-Z19C}
\textbf{Problem:} ESP32-S3 hat nur 2 Hardware-UARTs (UART0 für USB-Monitor, UART1 für PMS5003), aber 3 UART-Sensoren.

\textbf{Lösung:} EspSoftwareSerial (v8.2.0) für MH-Z19C bei 9600\,Baud auf GPIO4/5. Die niedrige Baudrate ermöglicht zuverlässigen Betrieb per Software-UART.

% ============================================================
\section{Software-Architektur}
% ============================================================

\subsection{Schichtenmodell}
Die Software folgt einem dreischichtigen Architekturmodell:

\begin{lstlisting}[language={},basicstyle=\ttfamily\footnotesize]
+----------------------------------------------------------+
|  APPLICATION LAYER (main.cpp)                            |
|  - State Machine (Active, Dimmed, Light Sleep)           |
|  - Event Handling (Button, Presence)                     |
|  - Timing Control, WiFi Clock                            |
+----------------------------------------------------------+
|  SERVICE LAYER                                           |
|  +----------------+ +----------------+ +---------------+ |
|  | Sensor-Module  | | UI-Module      | | Utilities     | |
|  | - aht_sgp.*    | | - ui_manager.* | | - sensor_     | |
|  | - mhz19c.*     | | - lvgl_driver.*| |   filter.*    | |
|  | - pms5003.*    | | - colors.h     | | - sensor_     | |
|  | - ld2410c.*    | | - display_     | |   history.*   | |
|  |                | |   config.h     | | - power_      | |
|  |                | |                | |   manager.*   | |
|  +----------------+ +----------------+ +---------------+ |
+----------------------------------------------------------+
|  DRIVER LAYER (Libraries)                                |
|  LovyanGFX 1.1.16 | LVGL 9.2.2 | Adafruit AHTX0/SGP40 |
|  MH-Z19 | PMS Library | ld2410 | EspSoftwareSerial      |
+----------------------------------------------------------+
\end{lstlisting}

\subsection{Verzeichnisstruktur}

\begin{lstlisting}[language={},basicstyle=\ttfamily\footnotesize]
inspectair/
  platformio.ini            # PlatformIO Konfiguration
  src/
    main.cpp                # Hauptprogramm, State Machine
    config/
      pins.h                # GPIO Pin-Definitionen
      sensor_types.h        # Sensor-Datenstrukturen
      colors.h              # Farben & Grenzwerte
      display_config.h      # Display-Konfiguration
      app_config.h          # Applikations-Konfiguration
    sensors/
      aht_sgp.cpp/.h        # AHT20 + SGP40 (I2C)
      mhz19c.cpp/.h         # MH-Z19C CO2 (SoftwareSerial)
      pms5003.cpp/.h        # PMS5003 Feinstaub (UART1)
      ld2410c.cpp/.h        # LD2410C Radar (UART2)
      sensor_filter.cpp/.h  # Gleitender Mittelwert
      sensor_history.cpp/.h # 24h Datenspeicher
    display/
      ui_manager.cpp/.h     # LVGL 9 Multi-Screen UI
      lvgl_driver.cpp/.h    # LovyanGFX <-> LVGL Brücke
      fonts/                # Custom Fonts mit Umlauten
    network/
      WifiClock.cpp/.h      # NTP-Zeitsynchronisation
      web_remote.cpp/.h     # Web-basierte Fernsteuerung
    power/
      power_manager.cpp/.h  # Dimming & Light Sleep
  lib/                      # Externe Libraries (PlatformIO)
\end{lstlisting}

\subsection{Datenstrukturen}

Die zentrale Datenstruktur \texttt{SensorReadings} (definiert in \texttt{sensor\_types.h}) aggregiert alle Sensordaten:

\begin{lstlisting}[language=C++]
typedef struct {
  AHT20_Data aht;       // Temperatur/Feuchte
  SGP40_Data sgp;       // VOC-Daten
  PMS5003_Data pms;     // Feinstaub (PM1.0/2.5/10)
  MHZ19C_Data mhz;      // CO2
  LD2410C_Data radar;   // Praesenz/Distanz
  uint32_t timestamp;   // Zeitstempel (millis)
} SensorReadings;
\end{lstlisting}

Jeder Sensor hat eine eigene Datenstruktur mit typisierten Feldern und Validierungsflags.

\subsection{Display-Architektur (LVGL 9)}

Die Display-Schicht besteht aus zwei Modulen:

\textbf{lvgl\_driver} (Low-Level): Initialisiert LovyanGFX für den ST7796S-Controller und verbindet es als LVGL-Display-Treiber. Nutzt Dual-Buffer (je 480$\times$32 Pixel) in PSRAM falls verfügbar. Stellt \texttt{lvgl\_init()}, \texttt{lvgl\_loop()} und \texttt{lvgl\_set\_brightness()} bereit.

\textbf{ui\_manager} (High-Level): Verwaltet fünf verschiedene Screen-Designs:

\begin{longtable}{l l p{7cm}}
\toprule
\rowcolor{tableheader}
\textbf{Screen} & \textbf{Enum} & \textbf{Beschreibung} \\
\midrule
Tree & \texttt{UI\_SCREEN\_TREE} & Baum-Animation als Startbildschirm \\
Overview & \texttt{UI\_SCREEN\_OVERVIEW} & Große AQI-Anzeige + 2 Kacheln (minimalistisch) \\
Detail & \texttt{UI\_SCREEN\_DETAIL} & Kleine AQI + 4 Kacheln (alle Werte) \\
Analog & \texttt{UI\_SCREEN\_ANALOG} & Analoge Cockpit-Instrumente \\
Bubble & \texttt{UI\_SCREEN\_BUBBLE} & Dynamische Kreise (Modern Bubbles) \\
\bottomrule
\end{longtable}

Custom Fonts (12/16/20/28/48\,pt) mit deutschen Umlauten werden für die Darstellung verwendet.

\subsection{State Machine}

Das System verwendet eine State Machine für das Power Management:

\begin{lstlisting}[language={},basicstyle=\ttfamily\footnotesize]
  +---------------+
  |    ACTIVE     |<---------------------+
  |  (100% BL)    |                      |
  +-------+-------+                      | Praesenz
          | 60s keine Praesenz           | erkannt
          v                              |
  +-------+-------+                      |
  |    DIMMED     |----------------------+
  |   (15% BL)    |
  +-------+-------+
          | 5min keine Praesenz
          v
  +-------+-------+
  |  LIGHT SLEEP  | <-- Radar weckt per Interrupt
  | (~0,8mA)      |     Aufwachzeit ~2ms
  +---------------+
\end{lstlisting}

Light Sleep wurde Deep Sleep vorgezogen, da die schnellere Aufwachzeit (${\sim}$2\,ms vs. mehrere Sekunden) eine sofortige Reaktion auf Präsenzerkennung ermöglicht und Sensordaten im RAM erhalten bleiben.

\subsection{Verwendete Libraries}

\begin{longtable}{p{3.5cm} p{2cm} p{6.5cm}}
\toprule
\rowcolor{tableheader}
\textbf{Library} & \textbf{Version} & \textbf{Verwendung} \\
\midrule
\endhead

LovyanGFX & 1.1.16 & Display-Treiber für ST7796S (SPI), Backlight-PWM \\
LVGL & 9.2.2 & UI-Framework, 5~Screen-Designs, Custom Fonts \\
Adafruit AHTX0 & latest & AHT20 Temperatur/Feuchte-Sensor (I\textsuperscript{2}C) \\
Adafruit SGP40 & latest & SGP40 VOC-Sensor mit VOC-Index-Algorithmus \\
MH-Z19 & 1.5.4 & CO\textsubscript{2}-Sensor Treiber \\
PMS Library & 1.1.0 & Feinstaub-Sensor Treiber (PM1.0/2.5/10) \\
ld2410 & 0.1.3 & Radar-Sensor Treiber (Präsenz/Bewegung) \\
EspSoftwareSerial & 8.2.0 & Software-UART für MH-Z19C (GPIO4/5) \\

\bottomrule
\end{longtable}

% ============================================================
\section*{Änderungshistorie}
% ============================================================

\begin{longtable}{l l l p{7cm}}
\toprule
\rowcolor{tableheader}
\textbf{Version} & \textbf{Datum} & \textbf{Autor} & \textbf{Änderung} \\
\midrule
1.0 & Januar 2026 & Team & Erstversion \\
2.0 & Februar 2026 & Team & LVGL 8.3$\rightarrow$9.2.2. LovyanGFX 1.1.16 integriert. 5~Screens statt 2~UI-Modi. Deep Sleep durch Light Sleep ersetzt. Pin-Belegung aktualisiert. Modulare Code-Struktur dokumentiert. SoftwareSerial für MH-Z19C ergänzt. \\
\bottomrule
\end{longtable}

\end{document}
